\subsection{Collaborative counting Bloom filter}\label{sec:ccbf}

Our system records complaints in a special data structure which we
call a \emph{collaborative counting Bloom filter}, or CCBF. This data structure
shares some of the same basic functionality as a counting Bloom filter~\cite{TON:FCAB00, PODC:Mitzenmacher01} or count-min sketch~\cite{JAL:CorMut05}, which is to insert elements and
compute the (approximate) frequency of a given element.

Our CCBF differs from a usual count-min sketch in that each update
operation is accompanied by a \emph{user id}, and each user can only perform
a single update for a given element. This can be thought of as a strict generalization
of the normal count-min sketch operations, where the latter may be simulated by our
CCBF by choosing a unique user id for each update.

The actual data structure for the CCBF is also far simpler than the 2D array
of integers used for a count-min sketch; instead, we store only a single length-$s$
bit vector $T$. As a result, our CCBF will have the following desirable properties:

\begin{itemize}
    \item The bit-length of $T$ scales linearly with the total
    number of insertions.
    \item Each witness operation (insertion) changes exactly one bit in the underlying
    bit vector from 0 to 1.
    \item The CCBF is \emph{item-oblivious}, meaning that after observing
    an interactive update protocol, the adversary learns which user id made the update,
    but not which item was updated.
\end{itemize}

The downside to our CCBF is a far lower accuracy of the count operation in general compared to count-min sketches.
However, we will show that, for careful parameter choices, the count operation is highly
accurate within a certain range, which is precisely what is needed for the
current application.

\subsubsection{CCBF Construction}

The CCBF consists of a single size-$s$ bit vector $T$ and two operations:

\newcommand{\Increment}{\ensuremath{\mathsf{Increment}}}
\newcommand{\TestCount}{\ensuremath{\mathsf{TestCount}}}

\begin{itemize}
    \item $\Increment(x, C)$: Increases the count by 1 for item $x$
    according to user id $C$.
    \item $\TestCount(x, t)$: Returns true if the number of increments performed so far for item $x$ is \emph{probably} greater than or equal to $t$.
\end{itemize}

Note that $\TestCount$ is probabilistic, in the sense that it may return false
when the actual count is greater than $t$, or true when the actual count is less
than $t$. Our construction guarantees the correctness probability is always at least
$\tfrac{1}{2}$, and our tail bounds below show the correctness probability quickly
goes towards $1$ when the actual count is much smaller or larger than $t$.

The performance and accuracy of the CCBF is governed by three integer parameters
$s$, $u$, and $v$, with $u,v\le s$,
which must be set at construction time. The first, $s$, is the fixed size of the
table $T$. Each user $i$ is randomly assigned a static set of exactly $u$ locations in
the $T$; i.e., a uniformly random subset of $\{0,1,2,\ldots,s-1\}$, which we
call the \emph{user set}.
Similarly, each possible item $x$ is assigned a random set of exactly $v$
bit vector locations, which we call the \emph{item set}.

The two CCBF operations can be implemented by a single server and any number of clients. The protocols are simple and straightforward, save for the calculation of the \emph{tipping point} $\tau$ which we present in the next subsection.

In these protocols, the size-$s$ bit vector $T$ is considered \emph{public} or \emph{world-readable}; it is known by all parties at all times. In reality, the server who actually stores $T$ may send it to the client periodically, or whenever a client initiates a \Increment{} or \TestCount{} protocol. However, the bit vector $T$ is only writable by the server.

The $\Increment(x,C)$ protocol, outlined in \cref{alg:increment}, involves the User attempting to set a single bit from 0 to 1 within the item set for $x$. However, the user is only allowed to write locations within their own user set. So, if there are no 0 bits in the intersection of these two index sets, the user instead changes any other arbitrary 0 bit in its own user set in order to maintain item obliviousness.

\floatname{algorithm}{Algorithm}

\begin{algorithm}[htbp!]
    \caption{$\Increment(x,C)$\label{alg:increment}}
\begin{enumerate}
    \item User and server separately compute the list of $u$ user locations for user $C$, $U_C \subseteq \{0,\ldots,s-1\}$.
    \item User computes list of $v$ item locations for item $x$,
   $V_x \subseteq \{0,\ldots,s-1\}$
   \item User checks each location in $U_C$ in the table $T$ to compute a list $S_C = \{i\in U_C \mid T[i] = 0\}$ of \emph{settable} locations for user $C$
   \item If $S_C = \emptyset$, then the user cannot proceed and calls \textbf{abort}.
   \item Else if $S_C \cap V_x \ne \emptyset$, user picks a uniformly random index $i\gets S_C \cap V_x$ and sends index $i$ to server.
   \item Else user picks a random index $i\gets S_C$ and sends index $i$ to server.
   \item Server checks that received index $i$ is in the user set $U_C$ and that $T[i] = 0$, then sets $T[i]$ to 1.
\end{enumerate}
\end{algorithm}

Since the bit vector $T$ is considered world-readable, the only communication here is the single index $i$ from client to server over an authenticated channel. In reality, to avoid race conditions, the server will actually send the table entry values $T[i]$ for all $i\in U_C$ to the user first and lock the state of the global bit vector $T$ until receiving the single index response back from the user.

The $\TestCount(x,t)$ protocol is not interactive as it only requires reading the entries of $T$. The precise computation of the \emph{tipping point} $\tau$ is detailed in the next section. Note that this computation depends only on the \emph{total} number of bits set in the bit vector $T$ as well as the parameters $s,u,v$; therefore the computation of $\tau$ is independent of the item $x$ and could for example be performed once by the server and saved without violating item obliviousness.

This protocol is detailed in \cref{alg:testcount}.

\begin{algorithm}[htp!]
    \caption{$\TestCount(x,t)$\label{alg:testcount}}
\begin{enumerate}
    \item Use parameters $s,u,v$ and current value of $m$ total number of bits set in $T$, to compute the tipping point $\tau$.
    \item Compute list of $v$ item locations for item $x$, $V_x \subseteq \{0,\ldots,s-1\}$
    \item Check how many bits of $T$ are set for indices in $V_x$. Return \textbf{true} if and only if this count is greater than or equal to $\tau$.
\end{enumerate}
\end{algorithm}

\subsubsection{Calculating the tipping point}

The key to correctness of the \TestCount{} protocol is a calculation of the \emph{tipping point} $\tau$, which is the expected number of 1 bits within any item set, if that item has been incremented $t$ times. We now derive an algorithm to compute this expected value exactly, in $O(tv)$ time and $O(v)$ space.

Let $s$ be the total size of the table $T$ and $m\le s$ be the total number of calls to \Increment{} so far. That is, $m$ equals the number of 1 bits in $T$.
Recall that $u,v \le s$ are the number of table entries per user and per item, respectively.

We first derive the probability that two subsets of the $s$ slots have given-size intersection.
Next we derive a recursive formula for $\tau$ using these intersection probabilities.
The nearest integer to $\tau$ can then be efficiently computed using a simple dynamic
programming strategy.

\paragraph{Intersection probabilities}

For the remainder, we use Knuth's notation $\falling{n}{k}$ to denote the \emph{falling factorial},
defined by
\[\falling{n}{k} = \frac{n!}{(n-k)!} = n\cdot(n-1)\cdot(n-2)\cdots(n-k+1).\]

\begin{lemma}\label{lem:intersect}
Let $k,a,b,s$ be non-negative integers with $k \le b \le a \le s$, and
suppose $S$ and $T$ are two subsets of a size-$s$ set with $|S|=a$ and $|T|=b$, each
chosen independently and uniformly over all subsets with those sizes. Then
\begin{equation}\label{eqn:intersect}
\Pr(|S \cap T| = k) = 
    \frac{\falling{a}{k} \cdot{} \falling{b}{k} \cdot{} \falling{(s-a)}{b-k}}
         {\falling{s}{b} \cdot{} k!}.
\end{equation}
\end{lemma}
\begin{proof}
The number of ways to choose $S$ and $T$ with a size-$k$ intersection, divided
by the total number of ways to choose two size-$a$ and size-$b$ sets, equals
\[\frac{\binom{s}{k} \cdot{} \binom{s-k}{a-k} \cdot{} \binom{s-a}{b-k}}
{\binom{s}{a}\cdot{}\binom{s}{b}}.\]
This simplifies to \eqref{eqn:intersect}.
\end{proof}

Because the numerator and denominator are each products of $b+k$ single-precision integers,
the value of \eqref{eqn:intersect} can be computed in $O(b)$ time to full accuracy in
machine floating-point precision.

Furthermore, equation \eqref{eqn:intersect} has the convenient property that, after altering any
value $a$, $b$, or $k$ by $\pm 1$, we can update the probability with only $O(1)$ additional
computation. So, for example, one can compute the probabilities for every $k\le b$ in the same
total time $O(b)$.

\paragraph{Recurrence for number of unfilled message slots}

Fix an arbitrary item $x$, and let $w\le v$ denote the number of 0 bits of $T$ within $x$'s item set. Let $k\le m$ denote the number of \Increment{} operations performed on item $x$ performed so far.

First, for convenience define $p_w$ to be the probability that an arbitrary user is able to
write to one of the $w$ remaining unfilled slots for the message. From \cref{lem:intersect}, we have
\begin{equation}
    p_w = 1 - \frac{\falling{(s-u)}{w}}{\falling{s}{w}},
    \label{eqn:pw}
\end{equation}
which can be computed in $O(w)$ time. In fact, we pre-compute \emph{all} possible values of $p_w$
with $0\le w\le v$ in $O(v)$ total time.

Now consider the random variable for the number of 0 bits within $x$'s item set after $k$ \Increment{}'s on $x$, if the item set originally had $w$ 0 bits.
Define $R_{w,k}$ to be the expected value of this random variable, which can be calculated recursively as follows.

If $w=0$, then the slots are all filled, and if $k=0$ then there are no more \Increment{}'s, so
the number of unfilled slots remains at $w$. Otherwise, the first \Increment{} will fill an additional slot
with probability $p_w$, leaving $w-1$ remaining unfilled slots, and otherwise will leave
$w$ remaining unfilled slots. This implies the following recurrence relation:

\[
R_{w,k} = \begin{cases}
    0,& w=0 \\
    w,& k=0 \\
    p_w R_{w-1,k-1} + (1 - p_w) R_{w,k-1},& w,k\ge 1
\end{cases}
\]

All values of $R_{w,t}$ with $0\le w \le v$ can be computed in $O(tv)$ time and $O(v)$ space,
using a straightforward dynamic programming strategy.

\paragraph{Computing the tipping point}

We now show how to compute the tipping point value $\tau$, which is the expected number of filled item slots
after $t$ \Increment{}s on that item, by summing the $R_{w,t}$ values over all possible
values of $w$ based on the number of \emph{other} calls to \Increment{} $m$.

To this end, define $q_w$ to be the probability that $w\le v$ slots for a given item are unfilled
after $m$ total calls to \Increment{} for other items. Because other calls to \Increment{} are for other
unrelated items, each one goes to a uniformly-random unfilled slot over the entire size-$s$ table $T$.
Therefore $q_w$ is the same as the probability of a size-$m$ set and a size-$v$ set having intersection size exactly $v-w$. From \cref{lem:intersect}, this is
\[q_w = 
    \frac{\falling{m}{v-w} \cdot{} \falling{v}{v-w} \cdot{} \falling{(s-m)}{w}}
         {\falling{s}{v} \cdot{} (v-w)!}.
\]
We can pre-compute all values of $q_w$ for $0\le w\le v$ in total time $O(v)$.

After pre-computing the values of $p_w$, $R_{w,t}$, and $q_w$, we can finally express
the tipping point $\tau$ as a linear combination
\begin{equation}\label{eqn:r}
    \tau = v - \sum_{w=0}^v q_w R_{w,t},
\end{equation}
rounded to the nearest integer.

In total, the computation requires $O(tv)$ time and $O(v)$ space.

\subsubsection{Tail Bounds}\label{sec:tailbounds}

\newcommand{\CminvOverk}{7.042652}
\newcommand{\CminukOvers}{0.5184846}
\newcommand{\Cplow}{0.956414}
\newcommand{\CmaxvOvers}{0.00386}
\newcommand{\CmaxuvOvers}{3.65151}
\newcommand{\Cphigh}{0.974876}
\newcommand{\CzeroUpper}{3.66567}
\newcommand{\CvLower}{371}
\newcommand{\CboxUpper}{0.00989}
\newcommand{\CminsOvern}{96}
\newcommand{\CmaxvOvert}{7.409}
\newcommand{\CmaxvOvertb}{8}
\newcommand{\Cfour}{1.0520553}
\newcommand{\Cminfp}{2.1}

Next, we prove lower and upper bounds on the probability of filling a single additional item slot during an \Increment{} operation, \cref{lem:lowerp,lem:upperp} respectively.
The proofs, which are intricate but not especially surprising, can be found in \cref{app:proofs}.

In order to make our scheme practically realizable, we state and prove explicit rather than asymptotic results, with all constants specified. These constants in themselves are not particularly meaningful; rather, they represent the tightest values which worked with our proof techniques and the parameter ranges we deemed reasonable for the application in mind.

\begin{restatable}{lemma}{lemmalowerp}\label{lem:lowerp}
    Let $x$ be an item such that at most
    $\tau$ of $x$'s item slots are filled.
    If the CCBF parameters $s,u,v$
    satisfy
    $v \ge \CminvOverk{} \tau$ and
    $u \ge \CminukOvers{} \tfrac{s}{\tau}$,
    then the probability that a call to $\Increment(x,C)$ fills in
    one more of $x$'s item slots is at least
    $\Cplow$.
\end{restatable}

\begin{restatable}{lemma}{lemmaupperp}\label{lem:upperp}
    Let $x$ be any item.
    If the CCBF parameters $s,u,v$ satisfy
    $\CvLower{} \le v \le \CmaxvOvers{} s$ and
    $u \le \CmaxuvOvers{} \tfrac{s}{v}$,
    then the probability that a call to $\Increment(x,C)$ fills in
    one more of $x$'s item slots is at most
    $\Cphigh$.
\end{restatable}

Now we use the probability upper bound to prove an upper bound on the tipping point $\tau$.

\begin{restatable}{lemma}{lemmauppertau}\label{lem:uppertau}
    Let $s,u,v$ be CCBF parameters that satisfy the conditions of \cref{lem:upperp},
    and suppose $m,t$ are integers such that
    $s \ge \CminsOvern{} m$ and
    $v \le \CmaxvOvert{} t$. Then the tipping point $\tau$, for threshold $t$
    and with $m$ total set bits in the table $T$, is at most
    $\Cfour{} t$.
\end{restatable}

We can now state our main theorems on the accuracy of the CCBF data structure. Consider a call to the predicate function $\TestCount(x,t)$, which attempts to determine whether the number of prior $\Increment{}$ calls with the same item $x$ is at least $t$. Our exact computation of the tipping point $r(t)$ shows that this function always returns the correct answer with at least 50\% probability. But of course, so would a random coin flip!

Let $k$ be the \emph{actual} number of calls to $\Increment(x,C)$ that have occurred. Then two kinds of errors can occur: a \emph{false positive} if $\TestCount(x,t)$ returns true but $k<t$, and a \emph{false negative} if $\TestCount(x,t)$ returns false when $k\ge t$. Intuitively, both errors occur with higher likelihood when the true count $k$ is close to $t$. Our main theorem captures and quantifies this intuition, saying that, ignoring low-order terms, $\TestCount$ is accurate to within a 10\% margin of error with high probability.

\begin{restatable}{theorem}{theoremfp}\label{thm:fp}
    Let $n$ be an upper bound on the total number of calls to $\Increment$, and $t$ be a desired threshold for $\TestCount$. Suppose the parameters $s,u,v$ for a CCBF data structure satisfy the conditions of \cref{lem:lowerp}, and furthermore that
    $v \le \CmaxvOvertb{} t$.
    If the actual number of calls to $\Increment(x,C)$ is at most
    $t - \Cminfp\sqrt{\lambda t}$, then the probability $\TestCount(x,t)$ gives
    a false positive is at most $2^{-\lambda}$.
\end{restatable}

\begin{restatable}{theorem}{theoremfn}\label{thm:fn}
    Let $n$ be an upper bound on the total number of calls to $\Increment$, and $t$ be a desired threshold for $\TestCount$. Suppose the parameters $s,u,v$ for a CCBF data structure satisfy the conditions 
    of \cref{lem:lowerp,lem:uppertau}.
    If the actual number of calls to $\Increment(x,C)$ is at least
    \begin{equation} \label{eqn:fn}
        1.1 t + .4 \lambda + .7\sqrt{\lambda t},
    \end{equation}
    then the probability $\TestCount(x,t)$ gives
    a false negative is at most $2^{-\lambda}$.
\end{restatable}

We can easily summarize the various conditions on the parameters as follows:
\begin{corollary}\label{cor:params}
    Let $n$ be a limit on the total number of calls to \Increment{}, and
    $t$ be a desired threshold satisfying 
    $50 \le t \le \tfrac{n}{20}$.
    Then by setting the parameters of a CCBF according to
    $s=96n$, $v=7.409t$, and $u=47.31\tfrac{n}{t}$, any call to
    $\TestCount(x,t)$ will satisfy the high accuracy assurances of
    \cref{thm:fp,thm:fn}.
\end{corollary}

%\todo{something about set of malicious users to strengthen first theorem}